\section{Methodology}
\label{sec:methodology}

Our methodology isolates feedback ordering from information volume through a matched-content experimental design. We describe the design rationale, conditions, and evaluation approach.

\subsection{Design Rationale}

Building on our observation that feedback ordering may independently affect repair quality, we design an experiment that controls the primary confound in prior studies: information volume. Rather than comparing ``static feedback'' vs ``execution feedback'' (which varies content), we compare ``static$\rightarrow$execution'' vs ``execution$\rightarrow$static'' (which varies only order).

\textbf{Key Design Principle:} Both conditions receive identical feedback content; only presentation order differs.

This requires:
\begin{enumerate}
    \item Generating both static and execution feedback for each problem
    \item Truncating each to a fixed token budget (500 tokens per type, 1000 total)
    \item Concatenating in different orders for each condition
    \item Ensuring byte-identical content across conditions
\end{enumerate}

\subsection{Experimental Conditions}

We define two matched conditions:

\textbf{Condition A (Static$\rightarrow$Execution):}
\begin{verbatim}
[Static Analysis Feedback: 500 tokens]
[Execution Feedback: 500 tokens]
\end{verbatim}

\textbf{Condition B (Execution$\rightarrow$Static):}
\begin{verbatim}
[Execution Feedback: 500 tokens]
[Static Analysis Feedback: 500 tokens]
\end{verbatim}

The feedback content is generated once per problem-iteration pair, then presented in both orders. Deterministic truncation (first 500 tokens) ensures identical content.

\subsection{Feedback Generation}

\textbf{Static Analysis:} We run Pylint and Mypy on generated code. Pylint catches style violations, unused variables, and potential bugs. Mypy catches type errors when type hints are present. Output is concatenated and truncated to 500 tokens.

\textbf{Execution Feedback:} We execute code against test cases in a sandboxed environment, capturing stdout/stderr, pass/fail status, and stack traces for failures. Output is truncated to 500 tokens, preserving traceback structure.

\subsection{Repair Loop}

Both conditions use identical generation/execution logic; only concatenation order differs. The loop terminates early if all tests pass or after 3 iterations.

\subsection{Evaluation Metrics}

\textbf{Primary Metric:} pass@1---fraction of problems where final code passes all tests.

\textbf{Relative Improvement:} (pass@1\_A - pass@1\_B) / pass@1\_B $\times$ 100\%

\textbf{Statistical Tests:} Bootstrap CI (10,000 resamples) for 95\% CI; McNemar's test for paired comparison.

\textbf{Mechanism Metrics:} $\Delta$Pass$_{1\rightarrow2}$ (change in pass rate from iteration 1 to 2); RegRate$_{1\rightarrow2}$ (P(pass@iter1 $\land$ fail@iter2)).
